% Supplementary derivation retained from the manuscript; not empirical data.
\section{Analytical Model}
\label{sec:analysis}

A simple model separates three effects that the experiments must distinguish:
capacity saved by sharing, the cost of repeated reads, and the opportunity
created by planning lead time. These equations describe conditional behavior;
they are not fitted to the measurements in Section~\ref{sec:evaluation}.

\subsection{Capacity and Access Cost}

Let $n$ branches each retain $S$ bytes, with a common immutable fraction
$\rho$ and otherwise distinct inputs. After producer handoff, private copies
require $nS$ bytes, while shared backing requires
$\rho S+n(1-\rho)S$. The payload saving is therefore $\rho(1-1/n)$.
Metadata, cache residency, retained source copies, and migration buffers must
be accounted for separately. A workload with no simultaneous authorized
overlap has no deduplication opportunity, regardless of branch count.

Suppose the pool and local DRAM provide bandwidths $B_p$ and $B_d$, and every
branch scans its input $r$ times. A copy policy stages all inputs locally; a
selective-sharing policy reads the common part from the pool and stages the
private part. With mapping cost $h$ per branch and compiler cost $c$, a simple
aggregate service-demand model is
\begin{align}
 D_{\mathrm{copy}} &= \frac{nS}{B_p}+\frac{nrS}{B_d}, \\
 D_{\mathrm{share}} &= \frac{nS(1-\rho+r\rho)}{B_p}
   +\frac{nrS(1-\rho)}{B_d}+nh+c.
\end{align}
It counts source reads and scan reads and excludes destination writes, useful
cache reuse, overlap, and queueing. Thus, it exposes a tradeoff rather than
predicting runtime.
Small objects may fail to amortize control overhead. Repeated uncached reads
may justify a local replica even when pooled backing saves capacity. Real
cache behavior can change this result, which is why the experiment varies
scan count and reports tuned copying separately from fresh allocation.

\subsection{Value of Lookahead}

For a speculative transfer taking $d=S/B_p$, let $L$ be its lead time,
$p$ the probability it is used, and $c_p$ the exposed decision cost. In an
otherwise idle path, expected exposed-stall saving is
\begin{equation}
 E[\mathrm{saving}]=p\min\{L,d\}-c_p,
\end{equation}
and unused traffic is $(1-p)S$. This assumes the hinted version is correct
when used and excludes interference. A late correct prediction may offer
little benefit; an early wrong prediction consumes bandwidth without
removing demand work. The controller must therefore consider accuracy,
lead time, and foreground pressure together.

The artifact's analytical script exports 412 deterministic parameter points
and the corresponding sensitivity plots. They remain separate from native
measurements. In particular, the empirical phase replay does not charge a
hypothetical compiler cost: it measures the native data path and reports
compiler overhead independently. Combining them requires knowing how much
compilation can overlap with the actual agent's planning interval.
